\documentclass[12pt, a4paper]{article}
\usepackage[utf8]{inputenc}
\usepackage[spanish]{babel}
\usepackage{graphicx}
\usepackage{float}
\usepackage{geometry}
\geometry{a4paper, margin=2.54cm}
\usepackage{hyperref}
\usepackage{grffile}
\usepackage{mathptmx} % Times New Roman
\usepackage{setspace}
\setstretch{1.5} % Interlineado APA 
\usepackage{xcolor}
\usepackage{titlesec}
\usepackage{array}
\usepackage{longtable}
\usepackage{booktabs}
\usepackage{listings}

% Configuración de títulos para APA
\titleformat{\section}{\normalfont\Large\bfseries}{\thesection}{1em}{}
\titleformat{\subsection}{\normalfont\large\bfseries}{\thesubsection}{1em}{}

% Configuración de listado de código (Java)
\definecolor{codegreen}{rgb}{0,0.6,0}
\definecolor{codegray}{rgb}{0.5,0.5,0.5}
\definecolor{codepurple}{rgb}{0.58,0,0.82}
\definecolor{backcolour}{rgb}{0.96,0.96,0.94}

\lstdefinestyle{javastyle}{
    backgroundcolor=\color{backcolour},   
    commentstyle=\color{codegreen},
    keywordstyle=\color{blue}\bfseries,
    numberstyle=\tiny\color{codegray},
    stringstyle=\color{codepurple},
    basicstyle=\ttfamily\footnotesize,
    breakatwhitespace=false,         
    breaklines=true,                 
    captionpos=b,                    
    keepspaces=true,                 
    numbers=left,                    
    numbersep=5pt,                  
    showspaces=false,                
    showstringspaces=false,
    showtabs=false,                  
    tabsize=4,
    language=Java
}
\lstset{style=javastyle}

\begin{document}

% Encabezado Institucional
\begin{center}
    \textbf{Universidad Tecnológica Israel} \\
    \textbf{Francisco Pizarro E4-142 y Marieta de Veintimilla} \\
    \textbf{Quito - Ecuador} \\
    \textbf{(593) 2 255-5741}
\end{center}
\noindent\rule{\textwidth}{1pt}

\vspace{0.5cm}
\noindent
\textbf{Facultad:} Ciencias de la Ingeniería \\
\textbf{Asignatura:} Fundamentos de la Especialidad II \\
\textbf{Docente:} Mg. Viviana Flores \\
\textbf{Estudiante:} Raul Fernando Cajiao Garces \\

\vspace{1cm}
\begin{center}
    \Large \textbf{Actividad Práctica: Auditoría y Refactorización de Código} \\
    \large \textbf{Aplicación de Principios de Diseño SOLID, DRY, KISS e ISO/IEC 25010}
\end{center}

\vspace{0.5cm}

\section{Introducción y Objetivos}
El desarrollo de software moderno no solo requiere que el código sea funcional, sino que también cumpla con estándares de calidad que faciliten su mantenimiento, escalabilidad y testabilidad a lo largo del tiempo. 

Esta actividad práctica tiene como objetivo auditar un componente denominado \texttt{GestorPedidos}, identificar violaciones de diseño clave a través del análisis de principios SOLID, DRY y KISS, y aplicar un proceso de refactorización estructural. Asimismo, se evalúan las mejoras del código utilizando las subcaracterísticas de la norma internacional \textbf{ISO/IEC 25010} (específicamente modularidad, reusabilidad, analizabilidad, modificabilidad y testabilidad), contrastando el estado original con la arquitectura propuesta.

\section{Enlaces del Entregable}
A continuación se detallan los enlaces correspondientes al desarrollo del laboratorio:
\begin{itemize}
    \item \textbf{Repositorio de GitHub:} \href{https://github.com/Raul-Fernando-Cajiao-Garces/DEBER1_FUNDAMETOS2}{https://github.com/Raul-Fernando-Cajiao-Garces/DEBER1\_FUNDAMETOS2}
    \item \textbf{Video Demostrativo (Explicación técnica):} \href{https://youtu.be/video-demostrativo-placeholder}{https://youtu.be/video-demostrativo-placeholder}
\end{itemize}

\newpage
\section{Código Original Auditado}
El siguiente listado muestra el código original en Java de la clase \texttt{GestorPedidos}, la cual contiene múltiples violaciones a buenas prácticas y principios de diseño de software.

\begin{lstlisting}[caption={Código Java Original}, label={lst:original_code}]
import java.util.*;
import java.io.*;
import java.sql.*;

public class GestorPedidos {
    private Connection conexionBD;

    public GestorPedidos() {
        try {
            this.conexionBD = DriverManager.getConnection(
                "jdbc:mysql://localhost:3306/tienda", "root", "admin123");
        } catch (SQLException e) {
            e.printStackTrace();
        }
    }

    public void procesarPedido(String nombreCliente, String emailCliente,
                               List<String> nombresProductos,
                               List<Double> preciosProductos,
                               List<Integer> cantidades,
                               String tipoCliente) {
        if (nombreCliente == null || nombreCliente.trim().isEmpty()) {
            System.out.println("Error: nombre de cliente invalido");
            return;
        }
        if (emailCliente == null || !emailCliente.contains("@")) {
            System.out.println("Error: email invalido");
            return;
        }

        double subtotal = 0;
        for (int i = 0; i < nombresProductos.size(); i++) {
            subtotal += preciosProductos.get(i) * cantidades.get(i);
        }

        double descuento = 0;
        if (tipoCliente.equals("VIP")) {
            descuento = subtotal * 0.20;
        } else if (tipoCliente.equals("FRECUENTE")) {
            descuento = subtotal * 0.10;
        } else if (tipoCliente.equals("REGULAR")) {
            descuento = subtotal * 0.05;
        } else if (tipoCliente.equals("NUEVO")) {
            descuento = 0;
        }

        double impuesto = (subtotal - descuento) * 0.12;
        double total = subtotal - descuento + impuesto;

        try {
            Statement stmt = conexionBD.createStatement();
            String sql = "INSERT INTO pedidos (cliente, total) VALUES ('"
                + nombreCliente + "', " + total + ")";
            stmt.executeUpdate(sql);
        } catch (SQLException e) {
            System.out.println("Error al guardar el pedido: " + e.getMessage());
        }

        try {
            FileWriter writer = new FileWriter("factura_" + nombreCliente + ".txt");
            writer.write("FACTURA\n");
            writer.write("Cliente: " + nombreCliente + "\n");
            for (int i = 0; i < nombresProductos.size(); i++) {
                writer.write(nombresProductos.get(i) + " x" + cantidades.get(i)
                    + " = $" + (preciosProductos.get(i) * cantidades.get(i)) + "\n");
            }
            writer.write("Subtotal: $" + subtotal + "\n");
            writer.write("Descuento: $" + descuento + "\n");
            writer.write("Impuesto: $" + impuesto + "\n");
            writer.write("TOTAL: $" + total + "\n");
            writer.close();
        } catch (IOException e) {
            System.out.println("Error al generar la factura: " + e.getMessage());
        }

        System.out.println("Enviando correo a " + emailCliente + "...");
        System.out.println("Asunto: Confirmacion de pedido");
        System.out.println("Cuerpo: Estimado " + nombreCliente + ", su pedido por $"
            + total + " ha sido procesado.");
        System.out.println("[LOG] Pedido procesado para " + nombreCliente
            + " - Total: " + total);
    }

    public void cancelarPedido(String nombreCliente, String emailCliente, int idPedido) {
        if (nombreCliente == null || nombreCliente.trim().isEmpty()) {
            System.out.println("Error: nombre de cliente invalido");
            return;
        }
        if (emailCliente == null || !emailCliente.contains("@")) {
            System.out.println("Error: email invalido");
            return;
        }

        try {
            Statement stmt = conexionBD.createStatement();
            String sql = "DELETE FROM pedidos WHERE id = " + idPedido;
            stmt.executeUpdate(sql);
        } catch (SQLException e) {
            System.out.println("Error al cancelar el pedido: " + e.getMessage());
        }

        System.out.println("Enviando correo a " + emailCliente + "...");
        System.out.println("Asunto: Cancelacion de pedido");
        System.out.println("Cuerpo: Estimado " + nombreCliente + ", su pedido #"
            + idPedido + " ha sido cancelado.");
    }
}
\end{lstlisting}

\newpage
\section{Identificación de Violaciones de Principios (Auditoría)}
A continuación, se detalla la auditoría del código original mediante la identificación de violaciones específicas a los principios SOLID, DRY y KISS:

\begin{longtable}{p{3.5cm} p{2.2cm} p{9.3cm}}
\caption{Tabla de Violaciones de Principios de Diseño} \label{tab:violaciones} \\
\toprule
\textbf{Principio} & \textbf{Líneas / Ubicación} & \textbf{Descripción Técnica de la Violación} \\
\midrule
\endfirsthead
\caption[]{Tabla de Violaciones de Principios de Diseño (continuación)} \\
\toprule
\textbf{Principio} & \textbf{Líneas / Ubicación} & \textbf{Descripción Técnica de la Violación} \\
\midrule
\endhead
\bottomrule
\endfoot

\textbf{SRP (Single Responsibility Principle)} & Clase completa & La clase tiene múltiples responsabilidades: conectarse a la base de datos, validar parámetros, calcular precios, registrar transacciones en SQL, generar facturas en texto plano, notificar a los clientes y registrar logs en consola. Cualquier cambio en estas áreas requiere modificar la misma clase. \\
\midrule

\textbf{OCP (Open/Closed Principle)} & Líneas 39 a 48 & La lógica de descuentos depende de un árbol de condicionales \texttt{if-else} basado en el parámetro \texttt{tipoCliente}. Si se añade un nuevo tipo de cliente (e.g., PREMIUM), se debe alterar la estructura lógica de la clase, exponiendo el código a errores de regresión. \\
\midrule

\textbf{DIP (Dependency Inversion Principle)} & Constructor y Métodos & Hay un alto acoplamiento a implementaciones concretas de infraestructura: se invoca directamente \texttt{DriverManager.getConnection}, \texttt{FileWriter} y las operaciones SQL vía \texttt{Statement} con concatenación de strings, impidiendo mockear o sustituir la base de datos o el sistema de archivos durante pruebas unitarias. \\
\midrule

\textbf{DRY (Don't Repeat Yourself)} & Líneas 19-26 y 87-94 / Líneas 79-84 y 99-103 & Las validaciones de cliente (\texttt{nombreCliente} y \texttt{emailCliente}) se repiten textualmente tanto en \texttt{procesarPedido} como en \texttt{cancelarPedido}. Adicionalmente, el formato de envío de correo se repite en ambos flujos. \\
\midrule

\textbf{KISS (Keep It Simple, Stupid)} & Métodos de negocio & Los métodos \texttt{procesarPedido} y \texttt{cancelarPedido} mezclan lógica de validación, computación matemática, persistencia de red e I/O de disco. Esto vuelve compleja la lectura del código, y propicia posibles fallos difíciles de depurar en un único bloque monolítico. \\
\midrule

\textbf{Seguridad (Extra)} & Líneas 50-52 y 90-91 & Vulnerabilidad de \textbf{Inyección SQL} debido a la concatenación directa de parámetros en sentencias ejecutadas por \texttt{Statement} en lugar de usar \texttt{PreparedStatement}. Además, las credenciales de la base de datos están quemadas (\textit{hardcoded}). \\
\bottomrule
\end{longtable}

\newpage
\section{Código Refactorizado Propuesto}
Para solventar las deficiencias identificadas, el código fue separado en módulos según su responsabilidad, inyectando dependencias por medio de abstracciones (interfaces) y aplicando el patrón de diseño \textbf{Strategy} para los descuentos.

\subsection{1. Abstracciones e Interfaces}
\begin{lstlisting}[caption={Abstracciones del Sistema}, label={lst:interfaces}]
package refactorizado.interfaces;
import refactorizado.model.Pedido;
import refactorizado.model.DetallePrecios;

public interface IPedidoRepository {
    void guardarPedido(String cliente, double total);
    void eliminarPedido(int idPedido);
}

public interface IFacturador {
    void generarFactura(Pedido pedido, DetallePrecios precios);
}

public interface INotificador {
    void enviarConfirmacion(String nombreCliente, String emailCliente, double total);
    void enviarCancelacion(String nombreCliente, String emailCliente, int idPedido);
}

public interface IDescuentoStrategy {
    double calcularDescuento(double subtotal);
}
\end{lstlisting}

\subsection{2. Patrón de Diseño Strategy (Descuentos)}
\begin{lstlisting}[caption={Estrategias de Descuento (SOLID OCP)}, label={lst:strategy}]
package refactorizado.strategy;
import refactorizado.interfaces.IDescuentoStrategy;
import java.util.HashMap;
import java.util.Map;

public class DescuentoVIP implements IDescuentoStrategy {
    public double calcularDescuento(double subtotal) { return subtotal * 0.20; }
}

public class DescuentoFrecuente implements IDescuentoStrategy {
    public double calcularDescuento(double subtotal) { return subtotal * 0.10; }
}

public class DescuentoRegular implements IDescuentoStrategy {
    public double calcularDescuento(double subtotal) { return subtotal * 0.05; }
}

public class DescuentoNuevo implements IDescuentoStrategy {
    public double calcularDescuento(double subtotal) { return 0.0; }
}

// Factoria de Estrategias para evitar condicionales
public class DescuentoFactory {
    private static final Map<String, IDescuentoStrategy> estrategias = new HashMap<>();

    static {
        estrategias.put("VIP", new DescuentoVIP());
        estrategias.put("FRECUENTE", new DescuentoFrecuente());
        estrategias.put("REGULAR", new DescuentoRegular());
        estrategias.put("NUEVO", new DescuentoNuevo());
    }

    public static IDescuentoStrategy obtenerEstrategia(String tipoCliente) {
        return estrategias.getOrDefault(tipoCliente, new DescuentoNuevo());
    }
}
\end{lstlisting}

\newpage
\subsection{3. Validadores y Lógica de Dominio}
\begin{lstlisting}[caption={Clases de Dominio, Validación y Cálculo de Precios}, label={lst:dominio}]
package refactorizado.model;
import java.util.List;

public class ValidadorCliente {
    public static void validar(String nombreCliente, String emailCliente) {
        if (nombreCliente == null || nombreCliente.trim().isEmpty()) {
            throw new IllegalArgumentException("Error: nombre de cliente invalido");
        }
        if (emailCliente == null || !emailCliente.contains("@")) {
            throw new IllegalArgumentException("Error: email invalido");
        }
    }
}

public class Pedido {
    private final String nombreCliente;
    private final String emailCliente;
    private final List<String> nombresProductos;
    private final List<Double> preciosProductos;
    private final List<Integer> cantidades;
    private final String tipoCliente;

    public Pedido(String nombreCliente, String emailCliente, List<String> nombresProductos,
                  List<Double> preciosProductos, List<Integer> cantidades, String tipoCliente) {
        ValidadorCliente.validar(nombreCliente, emailCliente);
        this.nombreCliente = nombreCliente;
        this.emailCliente = emailCliente;
        this.nombresProductos = nombresProductos;
        this.preciosProductos = preciosProductos;
        this.cantidades = cantidades;
        this.tipoCliente = tipoCliente;
    }

    public String getNombreCliente() { return nombreCliente; }
    public String getEmailCliente() { return emailCliente; }
    public List<String> getNombresProductos() { return nombresProductos; }
    public List<Double> getPreciosProductos() { return preciosProductos; }
    public List<Integer> getCantidades() { return cantidades; }
    public String getTipoCliente() { return tipoCliente; }

    public double calcularSubtotal() {
        double subtotal = 0;
        for (int i = 0; i < nombresProductos.size(); i++) {
            subtotal += preciosProductos.get(i) * cantidades.get(i);
        }
        return subtotal;
    }
}

public class DetallePrecios {
    public final double subtotal;
    public final double descuento;
    public final double impuesto;
    public final double total;

    public DetallePrecios(double subtotal, double descuento, double impuesto, double total) {
        this.subtotal = subtotal;
        this.descuento = descuento;
        this.impuesto = impuesto;
        this.total = total;
    }
}

public class CalculadoraPrecios {
    private static final double IVA_RATE = 0.12;

    public static DetallePrecios calcular(Pedido pedido) {
        double subtotal = pedido.calcularSubtotal();
        double descuento = DescuentoFactory.obtenerEstrategia(pedido.getTipoCliente())
                                           .calcularDescuento(subtotal);
        double impuesto = (subtotal - descuento) * IVA_RATE;
        double total = subtotal - descuento + impuesto;
        return new DetallePrecios(subtotal, descuento, impuesto, total);
    }
}
\end{lstlisting}

\newpage
\subsection{4. Implementaciones de Infraestructura}
\begin{lstlisting}[caption={Implementaciones Concretas}, label={lst:infraestructura}]
package refactorizado.infrastructure;
import refactorizado.interfaces.*;
import refactorizado.model.*;
import java.io.FileWriter;
import java.io.IOException;
import java.sql.Connection;
import java.sql.PreparedStatement;
import java.sql.SQLException;

public class PedidoRepositorySQL implements IPedidoRepository {
    private final Connection conexionBD;

    public PedidoRepositorySQL(Connection conexionBD) {
        this.conexionBD = conexionBD;
    }

    @Override
    public void guardarPedido(String cliente, double total) {
        String sql = "INSERT INTO pedidos (cliente, total) VALUES (?, ?)";
        try (PreparedStatement stmt = conexionBD.prepareStatement(sql)) {
            stmt.setString(1, cliente);
            stmt.setDouble(2, total);
            stmt.executeUpdate();
        } catch (SQLException e) {
            System.out.println("Error al guardar el pedido: " + e.getMessage());
        }
    }

    @Override
    public void eliminarPedido(int idPedido) {
        String sql = "DELETE FROM pedidos WHERE id = ?";
        try (PreparedStatement stmt = conexionBD.prepareStatement(sql)) {
            stmt.setInt(1, idPedido);
            stmt.executeUpdate();
        } catch (SQLException e) {
            System.out.println("Error al cancelar el pedido: " + e.getMessage());
        }
    }
}

public class FacturadorTexto implements IFacturador {
    @Override
    public void generarFactura(Pedido pedido, DetallePrecios precios) {
        try (FileWriter writer = new FileWriter("factura_" + pedido.getNombreCliente() + ".txt")) {
            writer.write("FACTURA\n");
            writer.write("Cliente: " + pedido.getNombreCliente() + "\n");
            for (int i = 0; i < pedido.getNombresProductos().size(); i++) {
                writer.write(pedido.getNombresProductos().get(i) + " x" + pedido.getCantidades().get(i)
                    + " = $" + (pedido.getPreciosProductos().get(i) * pedido.getCantidades().get(i)) + "\n");
            }
            writer.write("Subtotal: $" + precios.subtotal + "\n");
            writer.write("Descuento: $" + precios.descuento + "\n");
            writer.write("Impuesto: $" + precios.impuesto + "\n");
            writer.write("TOTAL: $" + precios.total + "\n");
        } catch (IOException e) {
            System.out.println("Error al generar la factura: " + e.getMessage());
        }
    }
}

public class NotificadorConsole implements INotificador {
    @Override
    public void enviarConfirmacion(String nombreCliente, String emailCliente, double total) {
        System.out.println("Enviando correo a " + emailCliente + "...");
        System.out.println("Asunto: Confirmacion de pedido");
        System.out.println("Cuerpo: Estimado " + nombreCliente + ", su pedido por $"
            + total + " ha sido procesado.");
        System.out.println("[LOG] Pedido procesado para " + nombreCliente + " - Total: " + total);
    }

    @Override
    public void enviarCancelacion(String nombreCliente, String emailCliente, int idPedido) {
        System.out.println("Enviando correo a " + emailCliente + "...");
        System.out.println("Asunto: Cancelacion de pedido");
        System.out.println("Cuerpo: Estimado " + nombreCliente + ", su pedido #"
            + idPedido + " ha sido cancelado.");
    }
}
\end{lstlisting}

\newpage
\subsection{5. Orquestador Refactorizado (Clase Principal)}
\begin{lstlisting}[caption={Clase GestorPedidos Refactorizada}, label={lst:gestor_refactored}]
package refactorizado;
import refactorizado.interfaces.*;
import refactorizado.model.*;

public class GestorPedidos {
    private final IPedidoRepository repository;
    private final IFacturador facturador;
    private final INotificador notificador;

    public GestorPedidos(IPedidoRepository repository, IFacturador facturador, INotificador notificador) {
        this.repository = repository;
        this.facturador = facturador;
        this.notificador = notificador;
    }

    public void procesarPedido(Pedido pedido) {
        DetallePrecios precios = CalculadoraPrecios.calcular(pedido);
        repository.guardarPedido(pedido.getNombreCliente(), precios.total);
        facturador.generarFactura(pedido, precios);
        notificador.enviarConfirmacion(pedido.getNombreCliente(), pedido.getEmailCliente(), precios.total);
    }

    public void cancelarPedido(String nombreCliente, String emailCliente, int idPedido) {
        ValidadorCliente.validar(nombreCliente, emailCliente);
        repository.eliminarPedido(idPedido);
        notificador.enviarCancelacion(nombreCliente, emailCliente, idPedido);
    }
}
\end{lstlisting}

\section{Justificación de la Refactorización (ISO/IEC 25010)}
El modelo de calidad de software ISO/IEC 25010 nos permite medir de forma estructurada los beneficios del proceso de refactorización sobre las características internas de mantenibilidad y modularidad.

\begin{longtable}{p{3.0cm} p{6.0cm} p{6.0cm}}
\caption{Comparación de Características de Calidad (Antes vs. Después)} \label{tab:iso25010} \\
\toprule
\textbf{Subcaracterística} & \textbf{Estado Anterior (Antes)} & \textbf{Estado Refactorizado (Después)} \\
\midrule
\endfirsthead
\caption[]{Comparación de Características de Calidad (continuación)} \\
\toprule
\textbf{Subcaracterística} & \textbf{Estado Anterior (Antes)} & \textbf{Estado Refactorizado (Después)} \\
\midrule
\endhead
\bottomrule
\endfoot

\textbf{Modularidad} & El código original residía en una única clase monolítica altamente acoplada. Modificar la persistencia afectaba a la lógica de negocio y viceversa. & Se divide en componentes independientes con interfaces claras: persistencia, facturación, y notificaciones. Los módulos pueden ser reemplazados sin impactar al resto. \\
\midrule

\textbf{Reusabilidad} & Lógica como el cálculo de descuentos o la generación de facturas estaba embebida dentro de un flujo lineal, imposible de reutilizar en otros servicios. & La calculadora de precios y las estrategias de descuento son reutilizables por cualquier otro orquestador o módulo del sistema. \\
\midrule

\textbf{Analizabilidad} & Comprender el código requería rastrear conexiones SQL manuales y concatenaciones junto a fórmulas aritméticas en más de 120 líneas de código secuencial. & El código es semánticamente autoexplicativo. Cada clase tiene menos de 40 líneas de código enfocadas en una única tarea estructurada. \\
\midrule

\textbf{Modificabilidad} & Agregar un tipo de descuento implicaba modificar la clase core (\texttt{GestorPedidos}), rompiendo potencialmente la persistencia o notificaciones. & Cumple con OCP. Para añadir un nuevo tipo de cliente solo se crea una clase que implemente \texttt{IDescuentoStrategy} y se registra en la fábrica, sin tocar el core. \\
\midrule

\textbf{Testabilidad} & Era imposible testear la lógica de negocio sin conectarse a una base de datos MySQL real y escribir archivos físicos en el sistema operativo. & Se permite la inyección de mocks para todas las interfaces auxiliares (\texttt{IPedidoRepository}, \texttt{INotificador}, etc.), posibilitando tests unitarios aislados del entorno externo. \\
\bottomrule
\end{longtable}

\newpage
\section{Historial del Proceso de Refactorización (Commits)}
Para demostrar la evolución e incrementalidad del proceso, se describe la secuencia planificada de confirmaciones (commits) realizadas en el repositorio de control de versiones Git:

\begin{enumerate}
    \item \textbf{Commit 1 - [Feat]:} Añadir código original de \texttt{GestorPedidos.java} como versión inicial sin modificaciones para establecer la línea base.
    \item \textbf{Commit 2 - [Refactor]:} Crear modelo de dominio \texttt{Pedido}, clase de utilidad \texttt{ValidadorCliente} y transferir la lógica de validación para eliminar código duplicado (DRY).
    \item \textbf{Commit 3 - [Pattern]:} Implementar interfaz \texttt{IDescuentoStrategy} y el patrón Strategy para los tipos de cliente VIP, Frecuente, Regular y Nuevo, eliminando la cadena condicional y cumpliendo con el principio OCP.
    \item \textbf{Commit 4 - [Arch]:} Definir abstracciones para los servicios de infraestructura (\texttt{IPedidoRepository}, \texttt{IFacturador}, \texttt{INotificador}) y migrar sus implementaciones para respetar los principios SRP y DIP.
    \item \textbf{Commit 5 - [Integrate]:} Actualizar la clase principal \texttt{GestorPedidos} para inyectar sus dependencias vía constructor y corregir vulnerabilidad de inyección SQL mediante \texttt{PreparedStatement}.
\end{enumerate}

\section{Conclusión}
La transición de un modelo transaccional procedural altamente acoplado a un diseño orientado a objetos basado en interfaces y patrones conductuales (\textit{Strategy} y \textit{Factory}) garantiza una estructura robusta. Esto permite a la organización tecnológica adaptarse a las demandas del mercado reduciendo el costo de cambio, eliminando errores colaterales comunes y mejorando significativamente la conformidad con la métrica ISO/IEC 25010 de mantenibilidad de software.

\end{document}
